\documentclass{article}
\usepackage[UTF8]{ctex}
\usepackage{amsmath}
\usepackage{amssymb}
\usepackage{graphicx}
\usepackage{geometry}
\geometry{a4paper, margin=2cm}

\newcommand{\problemname}[1]{\begin{center}\Large\bfseries #1\end{center}}
\usepackage{enumitem}

\begin{document}
\problemname{程序设计 T3}

蒜头君最爱吃草莓和芒果啦~ 这天幸福水果店正在进行多买多优惠活动，正好有蒜头君喜欢的水果！

水果店的优惠活动是这样的：如果购买的水果总重量在 $20$ 斤及以上，那么顾客可以选择自己购买的任一水果参与折扣活动，即单价打八折。

蒜头君购买了草莓若干，芒果若干，请你帮他计算一下，按照优惠规则，他最低要支付的价格是多少吧~

\subsection*{输入格式}
输入包含 $2$ 行：

\begin{itemize}
  \item 

第一行是两个数： $m(0<m\le100)$ ，表示购买的草莓的重量； $p_{1}(1<p_{1}\le100)$，表示草莓的单价；

  \item 

第二行是两个数： $n(0<n\le100)$ ，表示购买的芒果的重量； $p_{2}(1<p_{2}\le100)$，表示芒果的单价。

\end{itemize}
\subsection*{输出格式}
输出一个结果保留两位小数的数字，表示蒜头君最低要支付的金额。

\subsection*{样例输入}
\begin{verbatim}
10 8
10 9
\end{verbatim}

\subsection*{样例输出}
\begin{verbatim}
152.00
\end{verbatim}


\end{document}
